
\chapter{2016年江苏卷}

\section{填空题}

\begin{problem}
已知集合 \(A=\{-1,2,3,6\}\)，\(B=\{x\mid -2<x<3\}\)，则 \(A\cap B=\fillinblank{}\).

\end{problem}

\begin{answer}
\(\{-1,2\}\)
\end{answer}

\begin{solution}
因为 \(A=\{-1,2,3,6\}\)，\(B=\{x\mid -2<x<3\}\)，所以 \(A\cap B=\{-1,2\}\).
\end{solution}

\begin{problem}
复数 \(z=(1+2\i)(3-\i)\)，其中 \(\i\) 为虚数单位，则 \(z\) 的实部是\fillinblank{}.

\end{problem}

\begin{answer}
\(5\)
\end{answer}

\begin{solution}
由复数运算得 \(z=(1+2\i)(3-\i)=5+5\i\)，所以 \(z\) 的实部是 \(5\).
\end{solution}

\begin{problem}
在平面直角坐标系 \(xOy\) 中，双曲线 \(\dfrac{x^2}{7}-\dfrac{y^2}{3}=1\) 的焦距是\fillinblank{}.

\end{problem}

\begin{answer}
\(2\sqrt{10}\)
\end{answer}

\begin{solution}
由 \(a=\sqrt7\)，\(b=\sqrt3\)，得 \(c=\sqrt{a^2+b^2}=\sqrt{10}\)，故焦距为 \(2c=2\sqrt{10}\).
\end{solution}

\begin{problem}
已知一组数据 \(4.7,4.8,5.1,5.4,5.5\)，则该组数据的方差是\fillinblank{}.

\end{problem}

\begin{answer}
\(0.1\)
\end{answer}

\begin{solution}
这组数据的平均数为 \(5.1\)，所以方差
\(
\frac{1}{5}\bigl[(4.7-5.1)^2+(4.8-5.1)^2+(5.1-5.1)^2+(5.4-5.1)^2+(5.5-5.1)^2\bigr]=0.1
\).
\end{solution}

\begin{problem}
函数 \(y=\sqrt{3-2x-x^2}\) 的定义域是\fillinblank{}.

\end{problem}

\begin{answer}
\(\left[-3,1\right]\)
\end{answer}

\begin{solution}
由 \(3-2x-x^2\ge0\) 得 \(x^2+2x-3\le0\)，解得 \(x\in[-3,1]\).
\end{solution}

\begin{problem}
如图是一个算法的流程图，则输出的 \(a\) 的值是\fillinblank{}.

\bitmapfigure[max width=0.28\linewidth,max totalheight=4.8cm]{img/2016/jiangsu/q6_fig1.png}

\end{problem}

\begin{answer}
\(9\)
\end{answer}

\begin{solution}
模拟程序运行：初始 \(a=1\)，\(b=9\)，不满足 \(a>b\)，更新为 \(a=5\)，\(b=7\)；仍不满足，更新为 \(a=9\)，\(b=5\)；此时满足 \(a>b\)，输出 \(a=9\).
\end{solution}

\begin{problem}
将一颗质地均匀的骰子先后抛掷 \(2\) 次，则出现向上的点数之和小于 \(10\) 的概率是\fillinblank{}.

\end{problem}

\begin{answer}
\(\dfrac56\)
\end{answer}

\begin{solution}
基本事件总数为 \(6\times6=36\). 和不小于 \(10\) 的基本事件有 \((4,6)\)、\((6,4)\)、\((5,5)\)、\((5,6)\)、\((6,5)\)、\((6,6)\)，共 \(6\) 个，所以所求概率为 \(1-\dfrac{6}{36}=\dfrac56\).
\end{solution}

\begin{problem}
已知 \(\{a_n\}\) 是等差数列，\(S_n\) 是其前 \(n\) 项和，若 \(a_1+a_2^2=-3\)，\(S_5=10\)，则 \(a_9\) 的值是\fillinblank{}.

\end{problem}

\begin{answer}
\(20\)
\end{answer}

\begin{solution}
设公差为 \(d\)，则
\(a_1+(a_1+d)^2=-3\)，\(5a_1+\frac{5\cdot4}{2}d=10\).
\(
\begin{cases}
a_1+(a_1+d)^2=-3,\\
5a_1+10d=10,
\end{cases}
\)
解得 \(a_1=-4\)，\(d=3\)，故 \(a_9=a_1+8d=-4+24=20\).
\end{solution}

\begin{problem}
定义在区间 \([0,3\pi]\) 上的函数 \(y=\sin2x\) 的图象与 \(y=\cos x\) 的图象的交点个数是\fillinblank{}.

\end{problem}

\begin{answer}
\(7\)
\end{answer}

\begin{solution}
由
\(
\sin 2x=\cos x
\)
得
\(
2\sin x\cos x=\cos x
\)，
即
\(
\cos x\left( 2\sin x-1 \right)=0
\).

所以
\(
\cos x=0
\)
或
\(
\sin x=\dfrac12
\).

在区间 \([0,3\pi]\) 内，\(\cos x=0\) 给出
\(x=\dfrac{\pi}{2}\)，\(\dfrac{3\pi}{2}\)，\(\dfrac{5\pi}{2}\)，
共有 \(3\) 个解；
\(
\sin x=\dfrac12
\)
给出
\(x=\dfrac{\pi}{6}\)，\(\dfrac{5\pi}{6}\)，\(\dfrac{13\pi}{6}\)，\(\dfrac{17\pi}{6}\)，
共有 \(4\) 个解.

因此交点个数为 \(3+4=7\).
\end{solution}

\begin{problem}
如图，在平面直角坐标系 \(xOy\) 中，\(F\) 是椭圆 \(\dfrac{x^2}{a^2}+\dfrac{y^2}{b^2}=1\)（\(a>b>0\)） 的右焦点，直线 \(y=\dfrac{b}{2}\) 与椭圆交于 \(B,C\) 两点，且 \(\angle BFC=90^\circ\)，则该椭圆的离心率是\fillinblank{}.

\bitmapfigure[max width=0.45\linewidth,max totalheight=4.8cm]{img/2016/jiangsu/q10_fig1.png}

\end{problem}

\begin{answer}
\(\dfrac{\sqrt6}{3}\)
\end{answer}

\begin{solution}
设右焦点 \(F(c,0)\). 将 \(y=\dfrac b2\) 代入椭圆方程，得
\(
x=\pm a\sqrt{1-\frac{b^2}{4b^2}}=\pm\frac{\sqrt3}{2}a
\)，
故 \(B\left(-\dfrac{\sqrt3}{2}a,\dfrac b2\right)\)，\(C\left(\dfrac{\sqrt3}{2}a,\dfrac b2\right)\). 由 \(\angle BFC=90^\circ\) 得 \(k_{BF}k_{CF}=-1\)，化简得 \(b^2=3a^2-4c^2\). 又 \(b^2=a^2-c^2\)，所以 \(3c^2=2a^2\). 由 \(e=\dfrac ca\) 得 \(e=\dfrac{\sqrt6}{3}\).
\end{solution}

\begin{problem}
设 \(f(x)\) 是定义在 \(\R\) 上且周期为 \(2\) 的函数. 在区间 \([-1,1)\) 上，当 \(-1\le x<0\) 时，\(f(x)=x+a\)；当 \(0\le x<1\) 时，\(f(x)=\left|\dfrac25-x\right|\)，其中 \(a\in\R\). 若 \(f\!\left(-\dfrac52\right)=f\!\left(\dfrac92\right)\)，则 \(f(5a)\) 的值是\fillinblank{}.

\end{problem}

\begin{answer}
\(-\dfrac25\)
\end{answer}

\begin{solution}
由周期性，\(f\!\left(-\dfrac52\right)=f\!\left(-\dfrac12\right)=-\dfrac12+a\)，\(f\!\left(\dfrac92\right)=f\!\left(\dfrac12\right)=\left|\dfrac25-\dfrac12\right|=\dfrac{1}{10}\). 于是 \(a=\dfrac35\). 再由周期性得
\(
f(5a)=f(3)=f(-1)=-1+\frac35=-\frac25
\).
\end{solution}

\begin{problem}
已知实数 \(x,y\) 满足 \(x-2y+4\ge0\)、\(2x+y-2\ge0\)、\(3x-y-3\le0\)，则 \(x^2+y^2\) 的取值范围是\fillinblank{}.

\end{problem}

\begin{answer}
\(\left[\dfrac45,13\right]\)
\end{answer}

\begin{solution}
作出不等式组对应的平面区域，三个边界的交点分别为
\((0,2)\)，\((1,2)\)，\((2,3)\).
设
\(
z=x^2+y^2
\)，
则 \(z\) 表示区域内点到原点距离的平方.

由
\(0^2+2^2=4\)，\(1^2+2^2=5\)，\(2^2+3^2=13\)，
可得三个顶点到原点距离的平方分别为 \(4,5,13\)，其中 \(13\) 最大，所以点 \((2,3)\) 到原点的距离最大，从而
\(
z_{\max}=2^2+3^2=13
\).

又原点不在可行域内，可行域离原点最近的点在边界直线 \(2x+y-2=0\) 上，因此
\(
z_{\min}=\left(\frac{|{-2}|}{\sqrt{2^2+1^2}}\right)^2=\frac45
\).
所以取值范围为 \(\left[\dfrac45,13\right]\).
\end{solution}

\begin{problem}
如图，在 \(\triangle ABC\) 中，\(D\) 是 \(BC\) 的中点，\(E,F\) 是 \(AD\) 上的两个三等分点，\(\overrightarrow{BA}\cdot\overrightarrow{CA}=4\)，\(\overrightarrow{BF}\cdot\overrightarrow{CF}=-1\)，则 \(\overrightarrow{BE}\cdot\overrightarrow{CE}\) 的值是\fillinblank{}.

\bitmapfigure[max width=0.26\linewidth,max totalheight=4.8cm]{img/2016/jiangsu/q13_fig1.png}

\end{problem}

\begin{answer}
\(\dfrac78\)
\end{answer}

\begin{solution}
设 \(\overrightarrow{BD}=\bs d\)，\(\overrightarrow{DF}=\bs f\). 由 \(D\) 为 \(BC\) 中点，\(E,F\) 是 \(AD\) 上的两个三等分点，可得
\(\overrightarrow{BF}=\bs d+\bs f\)，\(\overrightarrow{CF}=-\bs d+\bs f\)，\(\overrightarrow{BA}=\bs d+3\bs f\)，\(\overrightarrow{CA}=-\bs d+3\bs f\)，
\(\overrightarrow{BE}=\bs d+2\bs f\)，\(\overrightarrow{CE}=-\bs d+2\bs f\).
由 \(\overrightarrow{BF}\cdot\overrightarrow{CF}=\bs f^2-\bs d^2=-1\)，\(\overrightarrow{BA}\cdot\overrightarrow{CA}=9\bs f^2-\bs d^2=4\)，解得 \(\bs f^2=\dfrac58\)，\(\bs d^2=\dfrac{13}{8}\). 于是
\(
\overrightarrow{BE}\cdot\overrightarrow{CE}=4\bs f^2-\bs d^2=\frac78
\).
\end{solution}

\begin{problem}
在锐角三角形 \(ABC\) 中，若 \(\sin A=2\sin B\sin C\)，则 \(\tan A\tan B\tan C\) 的最小值是\fillinblank{}.

\end{problem}

\begin{answer}
\(8\)
\end{answer}

\begin{solution}
由 \(\sin A=\sin(B+C)=\sin B\cos C+\cos B\sin C\) 及题设可得
\(
\sin B\cos C+\cos B\sin C=2\sin B\sin C
\).
因为 \(A,B,C\) 都是锐角，所以 \(\cos B,\cos C>0\)，可化为
\(
\tan B+\tan C=2\tan B\tan C
\).
又 \(\tan A=-\tan(B+C)=-\dfrac{\tan B+\tan C}{1-\tan B\tan C}\)，设 \(\tan B\tan C=t\)，则 \(t>1\)，
\(
\tan A\tan B\tan C=-\frac{2t^2}{1-t}=\frac{2t^2}{t-1}
\).
由 \(t>1\) 得该式最小值为 \(8\)，且当 \(t=2\) 时取到.
\end{solution}
\section{解答题}

\begin{problem}
在 \(\triangle ABC\) 中，\(AC=6\)，\(\cos B=\dfrac45\)，\(C=\dfrac{\pi}{4}\).

\begin{enumerate}
\item 求 \(AB\) 的长；
\item 求 \(\cos\!\left(A-\dfrac{\pi}{6}\right)\) 的值.
\end{enumerate}

\end{problem}

\begin{solution}
（1）由 \(\cos B=\dfrac45\) 得 \(\sin B=\dfrac35\). 由正弦定理，
\(
\frac{AB}{\sin C}=\frac{AC}{\sin B}
\)，
所以
\(
AB=\frac{6\cdot \sin(\pi/4)}{3/5}=5\sqrt2
\).

（2）由 \(\cos A=-\cos(B+C)=\sin B\sin C-\cos B\cos C=-\dfrac{\sqrt2}{10}\)，故
\(
\sin A=\sqrt{1-\cos^2A}=\frac{7\sqrt2}{10}
\).
于是
\(
\cos\!\left(A-\frac{\pi}{6}\right)=\frac{\sqrt3}{2}\cos A+\frac12\sin A=\frac{7\sqrt2-\sqrt6}{20}
\).
\end{solution}

\begin{problem}
如图，在直三棱柱 \(ABC-A_1B_1C_1\) 中，\(D,E\) 分别为 \(AB,BC\) 的中点，点 \(F\) 在侧棱 \(B_1B\) 上，且 \(B_1D\perp A_1F\)，\(A_1C_1\perp A_1B_1\). 求证：

\begin{enumerate}
\item 直线 \(DE\parallel\) 平面 \(A_1C_1F\)；
\item 平面 \(B_1DE\perp\) 平面 \(A_1C_1F\).
\end{enumerate}

\bitmapfigure[max width=0.32\linewidth,max totalheight=4.8cm]{img/2016/jiangsu/q16_fig1.png}

\end{problem}

\begin{solution}
（1）\(D,E\) 分别为 \(AB,BC\) 的中点，所以 \(DE\) 是 \(\triangle ABC\) 的中位线，故 \(DE\parallel AC\). 又因为 \(ABC-A_1B_1C_1\) 是棱柱，所以 \(AC\parallel A_1C_1\)，从而 \(DE\parallel A_1C_1\). 而 \(A_1C_1\subset\) 平面 \(A_1C_1F\)，故 \(DE\parallel\) 平面 \(A_1C_1F\).

（2）因为 \(ABC-A_1B_1C_1\) 是直棱柱，所以 \(AA_1\perp\) 平面 \(A_1B_1C_1\)，于是 \(AA_1\perp A_1C_1\). 又 \(A_1C_1\perp A_1B_1\)，故 \(A_1C_1\perp\) 平面 \(AA_1B_1B\). 由 \(DE\parallel A_1C_1\)，得 \(DE\perp\) 平面 \(AA_1B_1B\)，从而 \(DE\perp A_1F\). 另一方面，\(B_1D\perp A_1F\) 是题设条件，所以在平面 \(B_1DE\) 内，两条相交直线 \(B_1D\) 与 \(DE\) 都垂直于直线 \(A_1F\)，因此 \(A_1F\perp\) 平面 \(B_1DE\). 又 \(A_1F\subset\) 平面 \(A_1C_1F\)，所以平面 \(B_1DE\perp\) 平面 \(A_1C_1F\).
\end{solution}

\begin{problem}
现需要设计一个仓库，它由上下两部分组成，上部的形状是正四棱锥 \(P-A_1B_1C_1D_1\)，下部的形状是正四棱柱 \(ABCD-A_1B_1C_1D_1\)（如图所示），并要求正四棱柱的高 \(O_1O\) 是正四棱锥的高 \(PO_1\) 的 \(4\) 倍.

\begin{enumerate}
\item 若 \(AB=6\text{m}\)，\(PO_1=2\text{m}\)，则仓库的容积是多少？
\item 若正四棱锥的侧棱长为 \(6\text{m}\)，则当 \(PO_1\) 为多少时，仓库的容积最大？
\end{enumerate}

\bitmapfigure[max width=0.34\linewidth,max totalheight=4.8cm]{img/2016/jiangsu/q17_fig1.png}

\end{problem}

\begin{solution}
（1）由题意，正四棱柱的高 \(O_1O\) 是正四棱锥的高 \(PO_1\) 的 \(4\) 倍，故 \(PO_1=2\text{m}\) 时，\(O_1O=8\text{m}\). 仓库容积
\(
V=6^2\cdot8+\frac13\cdot6^2\cdot2=312\text{m}^3
\).

（2）设 \(PO_1=x\text{m}\)，则 \(O_1O=4x\text{m}\). 由侧棱长为 \(6\text{m}\)，得
\(A_1O_1=\sqrt{36-x^2}\)，\(A_1B_1=\sqrt2\,\sqrt{36-x^2}\).
于是仓库容积
\(V(x)=\left(\sqrt2\,\sqrt{36-x^2}\right)^2x+\frac13\left(\sqrt2\,\sqrt{36-x^2}\right)^2\cdot4x
=-\frac{26}{3}x^3+312x\)，\(0<x<6\).
所以
\(
V'(x)=-26x^2+312
\).
当 \(0<x<2\sqrt3\) 时，\(V'(x)>0\)；当 \(2\sqrt3<x<6\) 时，\(V'(x)<0\). 故 \(x=2\sqrt3\) 时取到最大值，即 \(PO_1=2\sqrt3\text{m}\) 时仓库的容积最大.
\end{solution}

\begin{problem}
如图，在平面直角坐标系 \(xOy\) 中，已知以 \(M\) 为圆心的圆 \(M:x^2+y^2-12x-14y+60=0\) 及其上一点 \(A(2,4)\).

\begin{enumerate}
\item 设圆 \(N\) 与 \(x\) 轴相切，与圆 \(M\) 外切，且圆心 \(N\) 在直线 \(x=6\) 上，求圆 \(N\) 的标准方程；
\item 设平行于 \(OA\) 的直线 \(l\) 与圆 \(M\) 相交于 \(B,C\) 两点，且 \(BC=OA\)，求直线 \(l\) 的方程；
\item 设点 \(T(t,0)\) 满足：存在圆 \(M\) 上的两点 \(P\) 和 \(Q\)，使得 \(\overrightarrow{TA}+\overrightarrow{TP}=\overrightarrow{TQ}\)，求实数 \(t\) 的取值范围.
\end{enumerate}

\bitmapfigure[max width=0.33\linewidth,max totalheight=4.8cm]{img/2016/jiangsu/q18_fig1.png}

\end{problem}

\begin{solution}
（1）圆 \(M\) 化为标准方程为 \((x-6)^2+(y-7)^2=25\)，其圆心为 \((6,7)\)，半径为 \(5\). 设圆 \(N\) 的圆心为 \((6,n)\)，且 \(n>0\). 因为圆 \(N\) 与 \(x\) 轴相切，所以其半径为 \(n\). 又与圆 \(M\) 外切，故
\(
|7-n|=n+5
\)，
解得 \(n=1\). 所以圆 \(N\) 的标准方程为 \((x-6)^2+(y-1)^2=1\).

（2）\(OA=\sqrt{2^2+4^2}=2\sqrt5\)，且直线 \(OA\) 的斜率为 \(2\)，故设 \(l:y=2x+b\). 圆心 \(M(6,7)\) 到直线 \(l\) 的距离为
\(
d=\frac{|12-7+b|}{\sqrt5}=\frac{|5+b|}{\sqrt5}
\).
圆 \(M\) 的半径为 \(5\)，所以
\(
BC=2\sqrt{25-d^2}
\).
由 \(BC=OA=2\sqrt5\)，得
\(
2\sqrt{25-\frac{(5+b)^2}{5}}=2\sqrt5
\)，
解得 \(b=5\) 或 \(b=-15\). 故 \(l\) 的方程为 \(y=2x+5\) 或 \(y=2x-15\).

（3）由 \(\overrightarrow{TA}+\overrightarrow{TP}=\overrightarrow{TQ}\) 得 \(\overrightarrow{TA}=\overrightarrow{PQ}\)，即 \(|TA|=|PQ|\). 圆 \(M\) 的半径为 \(5\)，故任意弦长 \(|PQ|\le10\). 又
\(
|TA|=\sqrt{(t-2)^2+4^2}
\).
故
\(
\sqrt{(t-2)^2+16}\le10
\)，
解得 \(t\in\left[2-2\sqrt{21},\,2+2\sqrt{21}\right]\).

对于任意 \(t\in\left[2-2\sqrt{21},\,2+2\sqrt{21}\right]\)，此时 \(|TA|\le10\). 只需要作直线 \(TA\) 的平行线，使圆心 \(M\) 到该直线的距离为
\(
\sqrt{25-\frac{|TA|^2}{4}}
\)，
则该直线必与圆 \(M\) 交于两点 \(P,Q\)，且
\(
|PQ|=2\sqrt{25-\left(25-\frac{|TA|^2}{4}\right)}=|TA|
\).
取 \(P,Q\) 的顺序使 \(\overrightarrow{PQ}\) 与 \(\overrightarrow{TA}\) 同向，则 \(\overrightarrow{PQ}=\overrightarrow{TA}\)，即 \(\overrightarrow{TA}+\overrightarrow{TP}=\overrightarrow{TQ}\).

因此实数 \(t\) 的取值范围为 \(t\in\left[2-2\sqrt{21},\,2+2\sqrt{21}\right]\).
\end{solution}

\begin{problem}
已知函数 \(f(x)=a^x+b^x\)（\(a>0,b>0,a\ne1,b\ne1\)）.

\begin{enumerate}
\item 设 \(a=2\)，\(b=\dfrac12\).

\circled{1} 求方程 \(f(x)=2\) 的根；

\circled{2} 若对于任意 \(x\in\R\)，不等式 \(f(2x)\ge mf(x)-6\) 恒成立，求实数 \(m\) 的最大值；
\item 若 \(0<a<1\)，\(b>1\)，函数 \(g(x)=f(x)-2\) 有且只有 \(1\) 个零点，求 \(ab\) 的值.
\end{enumerate}

\end{problem}

\begin{solution}
（1）设 \(a=2\)，\(b=\dfrac12\).

\circled{1} \(f(x)=2^x+2^{-x}\). 令 \(t=2^x\)（\(t>0\)），则
\(
t+\frac1t=2
\)，
解得 \(t=1\)，所以 \(x=0\).

\circled{2} 记 \(t=2^x+2^{-x}\)（\(t\ge2\)）. 原不等式化为
\(t^2-mt+4\ge0\)，\(t\ge2\).
若 \(m>4\)，则该二次函数在 \(t=\dfrac m2\) 处取最小值，无法恒成立；若 \(m\le4\)，则在 \([2,+\infty)\) 上单调递增，最小值取于 \(t=2\)，此时
\(
8-2m\ge0
\)，
即 \(m\le4\). 故 \(m\) 的最大值为 \(4\).

（2）设
\(
g(x)=a^x+b^x-2
\).
因为 \(0<a<1\)，\(b>1\)，所以 \(g(x)\) 有唯一极小值点 \(x_0\). 若 \(g(x)\) 只有一个零点，则极小值必为 \(0\)，从而 \(g(x_0)=0\). 由 \(g(0)=a^0+b^0-2=0\)，知唯一零点只能是 \(x_0=0\)，故
\(
a^0+b^0-2=0
\)
并结合极值点条件可得 \(\ln a+\ln b=0\)，即 \(ab=1\).
\end{solution}



\begin{problem}
如图，在 \(\triangle ABC\) 中，\(\angle ABC=90^\circ\)，\(BD\perp AC\)，\(D\) 为垂足，\(E\) 为 \(BC\) 的中点，求证：\(\angle EDC=\angle ABD\).

\bitmapfigure[max width=0.28\linewidth,max totalheight=4.8cm]{img/2016/jiangsu/q21_fig1.png}

\end{problem}

\begin{solution}
因为 \(BD\perp AC\)，所以 \(\angle BDC=90^\circ\).

因为 \(E\) 为 \(BC\) 的中点，所以在 \(\triangle BDC\) 中，\(DE=CE\)，则 \(\angle EDC=\angle C\).

由 \(\angle BDC=90^\circ\)，可得 \(\angle C+\angle DBC=90^\circ\).

由 \(\angle ABC=90^\circ\)，可得 \(\angle ABD+\angle DBC=90^\circ\).

因此 \(\angle ABD=\angle C\)，而 \(\angle EDC=\angle C\)，

所以 \(\angle EDC=\angle ABD\).
\end{solution}

\begin{problem}
已知矩阵 \(A=\begin{bmatrix}1&2\\0&-2\end{bmatrix}\)，矩阵 \(B\) 的逆矩阵 \(B^{-1}=\begin{bmatrix}1&-\dfrac12\\0&2\end{bmatrix}\)，求矩阵 \(AB\).

\end{problem}

\begin{solution}
由 \(B^{-1}=\begin{bmatrix}1&-\dfrac12\\0&2\end{bmatrix}\)，得
\(
B=(B^{-1})^{-1}=\begin{bmatrix}1&\dfrac14\\0&\dfrac12\end{bmatrix}
\).
因此
\(
AB=\begin{bmatrix}1&2\\0&-2\end{bmatrix}
\begin{bmatrix}1&\dfrac14\\0&\dfrac12\end{bmatrix}
=\begin{bmatrix}1&\dfrac54\\0&-1\end{bmatrix}
\).
\end{solution}

\begin{problem}
在平面直角坐标系 \(xOy\) 中，已知直线 \(l\) 的参数方程为 \(\left\{\begin{array}{l}x=1+\dfrac12 t,\\ y=\dfrac{\sqrt3}{2}t\end{array}\right.\)（\(t\) 为参数），椭圆 \(C\) 的参数方程为 \(\left\{\begin{array}{l}x=\cos\theta,\\ y=2\sin\theta\end{array}\right.\)（\(\theta\) 为参数），设直线 \(l\) 与椭圆 \(C\) 相交于 \(A,B\) 两点，求线段 \(AB\) 的长.

\end{problem}

\begin{solution}
由直线参数方程得
\(x=1+\frac12t\)，\(y=\frac{\sqrt3}{2}t\).
消去 \(t\) 得
\(
\sqrt3x-y-\sqrt3=0
\).
再由椭圆参数方程
\(x=\cos\theta\)，\(y=2\sin\theta\)，
可得椭圆普通方程为 \(x^2+\dfrac{y^2}{4}=1\). 联立
\(
\begin{cases}
\sqrt3x-y-\sqrt3=0,\\
x^2+\dfrac{y^2}{4}=1,
\end{cases}
\)
解得交点为 \((1,0)\) 和 \(\left(-\dfrac17,-\dfrac{8\sqrt3}{7}\right)\). 故
\(
AB=\sqrt{\left(1+\frac17\right)^2+\left(0+\frac{8\sqrt3}{7}\right)^2}=\frac{16}{7}
\).
\end{solution}

\begin{problem}
设 \(a>0\)，\(|x-1|<\dfrac a3\)，\(|y-2|<\dfrac a3\)，求证：\(|2x+y-4|<a\).

\end{problem}

\begin{solution}
由 \(a>0\)，\(|x-1|<\dfrac a3\)，\(|y-2|<\dfrac a3\)，得
\(
|2x+y-4|=\left|2(x-1)+(y-2)\right|
\le2|x-1|+|y-2|
<\frac{2a}{3}+\frac a3=a
\).
故 \(|2x+y-4|<a\) 成立.
\end{solution}

\begin{problem}
如图，在平面直角坐标系 \(xOy\) 中，已知直线 \(l:x-y-2=0\)，抛物线 \(C:y^2=2px\)（\(p>0\)）.

\begin{enumerate}
\item 若直线 \(l\) 过抛物线 \(C\) 的焦点，求抛物线 \(C\) 的方程；
\item 已知抛物线 \(C\) 上存在关于直线 \(l\) 对称的相异两点 \(P\) 和 \(Q\).

\circled{1} 求证：线段 \(PQ\) 的中点坐标为 \((2-p,-p)\)；

\circled{2} 求 \(p\) 的取值范围.
\end{enumerate}

\bitmapfigure[max width=0.34\linewidth,max totalheight=4.8cm]{img/2016/jiangsu/q25_fig1.png}

\end{problem}

\begin{solution}
（1）抛物线 \(y^2=2px\) 的焦点为 \(\left(\dfrac p2,0\right)\). 由于直线 \(l:x-y-2=0\) 过焦点，故 \(\dfrac p2-2=0\)，得 \(p=4\)，所以抛物线方程为 \(y^2=8x\).

（2）设 \(P(x_1,y_1)\)，\(Q(x_2,y_2)\).

\circled{1} 因为直线 \(l:x-y-2=0\) 的斜率为 \(1\)，而 \(P,Q\) 关于直线 \(l\) 对称，所以直线 \(PQ\perp l\)，从而
\(
k_{PQ}=-1
\).
又因为 \(P,Q\) 在抛物线 \(C\) 上，所以
\(x_1=\frac{y_1^2}{2p}\)，\(x_2=\frac{y_2^2}{2p}\).
所以
\(
k_{PQ}=\frac{y_1-y_2}{x_1-x_2}
=\frac{y_1-y_2}{\frac{y_1^2-y_2^2}{2p}}
=\frac{2p}{y_1+y_2}
\).
于是
\(
\frac{2p}{y_1+y_2}=-1
\)，
即 \(y_1+y_2=-2p\)，从而
\(
\frac{y_1+y_2}{2}=-p
\).
又 \(PQ\) 的中点在直线 \(l\) 上，故
\(
\frac{x_1+x_2}{2}-\frac{y_1+y_2}{2}-2=0
\)，
于是
\(
\frac{x_1+x_2}{2}=2-p
\)，
所以线段 \(PQ\) 的中点坐标为 \((2-p,-p)\).

\circled{2} 由 \circled{1} 知
\(\frac{y_1+y_2}{2}=-p\)，\(\frac{x_1+x_2}{2}=2-p\).
又
\(x_1=\frac{y_1^2}{2p}\)，\(x_2=\frac{y_2^2}{2p}\)，
所以
\(
\frac{y_1^2+y_2^2}{4p}=2-p
\)，
即
\(
y_1^2+y_2^2=8p-4p^2
\).
从而
\(
(y_1+y_2)^2=y_1^2+y_2^2+2y_1y_2
\)，
即
\(
4p^2=8p-4p^2+2y_1y_2
\)，
故
\(
y_1y_2=4p^2-4p
\).
于是 \(y_1,y_2\) 是方程
\(
y^2+2py+4p^2-4p=0
\)
的两个不相等的实根. 故判别式大于 \(0\)，即
\(
(2p)^2-4(4p^2-4p)>0
\)，
解得 \(0<p<\dfrac43\).
\end{solution}

\begin{problem}

\begin{enumerate}
\item 求 \(7\mathrm C_6^3-4\mathrm C_7^4\) 的值；
\item 设 \(m,n\in\N^*\)，\(n\ge m\)，求证：\((m+1)\mathrm C_m^m+(m+2)\mathrm C_{m+1}^m+(m+3)\mathrm C_{m+2}^m+\cdots+n\mathrm C_{n-1}^m+(n+1)\mathrm C_n^m=(m+1)\mathrm C_{n+2}^{m+2}\).
\end{enumerate}

\end{problem}

\begin{solution}
（1）
\(
7\mathrm C_6^3-4\mathrm C_7^4
=7\cdot\frac{6\cdot5\cdot4}{3\cdot2\cdot1}-4\cdot\frac{7\cdot6\cdot5\cdot4}{4\cdot3\cdot2\cdot1}
=7\cdot20-4\cdot35=0
\).

（2）对任意 \(m\in\N^*\).

当 \(n=m\) 时，左边等于 \((m+1)\mathrm C_m^m=m+1\)，右边等于 \((m+1)\mathrm C_{m+2}^{m+2}=m+1\)，等式成立.

假设当 \(n=k\)（\(k\ge m\)） 时命题成立，即
\(
(m+1)\mathrm C_m^m+(m+2)\mathrm C_{m+1}^m+\cdots+k\mathrm C_{k-1}^m+(k+1)\mathrm C_k^m
=(m+1)\mathrm C_{k+2}^{m+2}
\).
则当 \(n=k+1\) 时，左边
\(
=(m+1)\mathrm C_{k+2}^{m+2}+(k+2)\mathrm C_{k+1}^m
\).
由组合数恒等式可得
\(
(m+1)\mathrm C_{k+2}^{m+2}+(k+2)\mathrm C_{k+1}^m
=(m+1)\mathrm C_{k+3}^{m+2}
\)，
故 \(n=k+1\) 时等式也成立.

综上，对任意 \(m,n\in\N^*\) 且 \(n\ge m\)，有
\(
(m+1)\mathrm C_m^m+(m+2)\mathrm C_{m+1}^m+(m+3)\mathrm C_{m+2}^m+\cdots+n\mathrm C_{n-1}^m+(n+1)\mathrm C_n^m
=(m+1)\mathrm C_{n+2}^{m+2}
\).
\end{solution}
